\begin{abstract}
High-dimensional vector indices must balance construction time, memory, and query throughput. Two-bit comparisons reduce computation and storage, but preserving useful graph connections and search candidates requires scores suited to each stage. We present QuIVer, a graph index that uses a training-free sign--magnitude representation for construction and navigation while reserving original vectors for bounded final reranking. Construction discovers candidates with Hamming scores, orders them with an integer proxy, and defers reverse-edge pruning across parallel insertions. Search uses reconstruction norms and bitwise inner products to order both the frontier and retained candidates, without query calibration or an exact anchor. Experiments cover 21 datasets and compare 30 external baseline modes on three core workloads. QuIVer's validation-selected 95\% configurations lie on the measured Pareto frontiers of achieved recall, throughput, index preparation time, and memory, whether memory counts navigation state alone or also the original vectors. On Cohere-5M, QuIVer reaches 95.90\% Recall@10 at 3,075 single-thread queries/s, $11.74\times$ the throughput of Lucene BBQ with FP32 reranking at 95.83\% recall. Against SymphonyQG at 95.87\%, it retains 81\% of single-thread and 97\% of 32-thread throughput while reducing index preparation time by 54\%, total logical memory by 47\%, and navigation storage by 87\%. An SSD mode with INT8 reranking serves Cohere-5M at about 5 GiB of steady-state process RSS. Controlled crossovers and candidate-coverage diagnostics connect these results to construction and retention quality. Across 21 datasets, high recall holds for most text and image embeddings, while non-negative descriptors, cross-modal queries, and random data expose the limits of the representation.
\end{abstract}
